%% ============================================================================
\section{Feature Analysis}
\label{sec:features}
%% ============================================================================

This section characterises what the \resNcandAe~AE and \resNcandUs~US features respond to. We examine the operating-condition dependence of the features---distinguishing marginal from conditional correlations---and their low-dimensional structure.

\subsection{Operating-condition dependence}
\label{subsec:confound}

The raw correlation between signal features and $\kappa$ is partially driven by the shared dependence on rotational speed and temperature. Higher RPM reduces the required viscosity $\nureq$ (Eq.~\eqref{eq:nu1}) and therefore increases $\kappa$, while simultaneously increasing signal energy through more frequent contact events per second. Higher temperature reduces the actual viscosity
$\nuop$ and therefore reduces $\kappa$, while also affecting material damping and
structural dynamics.

As a consequence, amplitude features such as RMS are partially acting as indirect
proxies for rotational speed rather than as direct indicators of lubrication film
adequacy~\cite{hou_insitu_2025, zhang_coupled_2026, miettinen_analysis_2001}.
This confounding is inherent to any single-operating-point sensor system without an explicit operating-condition normalisation step. Because $\kappa$ is fully determined by the operating point, a model given measured RPM and temperature directly would predict $\kappa$ perfectly by construction; the contribution of the acoustic features is to recover this information without tachometer or thermometer access.

Marginally, the retained features correlate with $\kappa$ chiefly through this
operating-point dependence. Rotational speed is the dominant axis---it has a Spearman correlation of $\rho = \resRhoRpmKappa$ with $\kappa$, against $\rho = \resRhoTempKappa$ for temperature---and a LightGBM regressor on the retained AE features recovers the operating point itself almost completely (RPM at $R^2 = \resProxyRpmRsq$, temperature at $R^2 = \resProxyTempRsq$). Passing these predicted conditions through Eqs.~\eqref{eq:kappa}--\eqref{eq:nu1} yields a two-stage $\kappa$ estimate ($R^2 = \resTwoStageRsq$) that nearly equals the direct features-to-$\kappa$ model ($R^2 = \resHOrsqLgbAe$), so virtually all of the marginal skill is attributable to acoustic operating-point proxying rather
than to direct film sensing. This marginal view, however, is dominated by the
speed sweep and conflates the two $\kappa$ channels. We therefore condition on
each operating variable in turn, using the measured operating point: the
within-RPM-step correlation with temperature (sweeps grouped into staircase steps
of fixed speed) isolates the temperature response, while the within-temperature-block correlation with speed (sweeps grouped into
2~\si{\celsius} blocks of measured temperature, within which the temperature
standard deviation stays below 1~\si{\celsius}) isolates the speed response.
Figure~\ref{fig:marg_cond} contrasts the marginal and conditional structure for
every retained feature of both sensors. The speed axis is essentially unchanged
by conditioning---the features genuinely track rotational speed---whereas the
temperature axis spreads outward once speed is removed, uncovering a
lubrication-state sensitivity that the marginal view masks; the strongest AE
feature, \resWithinStepFeat, reaches a within-step correlation of $\rho = \resWithinStepRhoHC$ with temperature. The features are thus acoustic
operating-point soft sensors that, beyond the proxy share, retain a genuine
residual sensitivity to the lubrication state.

The two sensors occupy distinct regions of this conditional plane. AE features
cluster on the speed axis, whereas the passive-US channel resolves the temperature dependence more strongly and with mixed sign across its sub-bands
(Figure~\ref{fig:marg_cond}, right). AE and US therefore capture the lubrication
state through partly distinct channels rather than redundantly---the
complementarity their fusion later exploits.

\begin{figure}[H]
\centering
\includegraphics[width=\textwidth]{marginal_vs_conditional.png}
\caption{Marginal versus conditional rank-correlation structure of the retained
AE (circles) and US (squares) features. Each point is one feature; the axes are
the Spearman correlation with measured temperature ($x$) and rotational speed
($y$). \emph{Left:} marginal correlations over all sweeps, dominated by the speed
sweep. \emph{Right:} partition-based conditional correlations---within-RPM-step
for the temperature axis and within-temperature-block for the speed axis---with
whiskers showing the inter-quartile range across partitions. Conditioning leaves
the speed axis essentially unchanged but spreads the features along the
temperature axis, revealing lubrication-state sensitivity masked in the marginal
view. AE features track speed; US features lean toward temperature.}
\label{fig:marg_cond}
\end{figure}

\subsection{Low-dimensional structure and regime separability}
\label{subsec:pca}

Figure~\ref{fig:pca} shows the first two principal components of the retained AE and US feature sets coloured by $\kappa$ value. In the AE projection the $\kappa$ value varies smoothly across the plane rather than forming disjoint clusters: high-$\kappa$ sweeps (well-separated films) concentrate in the upper-left lobe while low-$\kappa$ sweeps fan out along a tail extending to large positive PC1, so the first two components---together capturing roughly two-thirds of the AE feature variance---already arrange the data along a continuous lubrication gradient. The US projection shows the same ordering more weakly: a single dense core carries a shallower $\kappa$ gradient and a detached low-$\kappa$ cluster sits apart from it, consistent with the lower predictive power of the US channel (Section~\ref{sec:modelling}). The continuity of the AE gradient, rather than sharp regime boundaries, mirrors the continuous nature of
$\kappa$ and indicates that the retained features encode a graded film-thickness signal rather than a small number of discrete operating states.

\begin{figure}[H]
\centering
\includegraphics[width=0.9\textwidth]{pca_kappa_regimes.png}
\caption{PCA projection of the retained AE and
US feature sets, coloured by lubrication $\kappa$ value.}
\label{fig:pca}
\end{figure}
